% ===================================================================
%  नकसा नं. ५ — घर आ ओकर बनावट।
%
%  Traduction, faite en 2026, en bhojpouri d'aujourd'hui, sur l'ido
%  de texto/io. Regles de la colonne : voir 10-nakasa-01.tex. Aucun
%  alinea n'est numerote par l'auteur : tout le tableau porte le
%  numero d'alinea « 01 », chaque replique formant un rang. 69 blocs,
%  142 renvois, TRENTE-SIX paires a retourner — le plus gros tableau
%  du livret. UNE note d'auteur, t05-noto-1, posee entre les deux
%  moities de la replique 25 ; la traduction ecrit UN bloc et pose la
%  note apres.
%
%  « मिस्त्री » LE MACON (108) EST PORTUGAIS, ET C'EST LE PREMIER DES
%  TROIS QUI NE SOIT PAS UN OBJET. La colonne gujaratie avait releve
%  મિસ્ત્રી parmi ses onze mots portugais — mestre, le maitre d'un
%  metier. Cette colonne-ci tenait deja मेज la table (tableau 1) et
%  पादरी le cure (tableau 3) ; celui-ci est le troisieme, et il ne
%  nomme ni un meuble ni un homme d'eglise mais UN METIER. Tout le
%  batiment de la plaine s'appelle par lui : राजमिस्त्री le macon,
%  काठ-मिस्त्री le menuisier, बिजली मिस्त्री l'electricien. Un mot
%  qui devient un suffixe de metier est un mot qui a cesse d'etre
%  etranger.
%
%  ET L'ECHAFAUDAGE PARLE CINQ LANGUES. Sur le seul chantier de ce
%  tableau : मिस्त्री le macon est portugais ; ठेकेदार l'entrepreneur
%  (77) est persan par le ठेका, le marche passe ; साहुल le fil a
%  plomb (109) est l'arabe شاقول ; मजूर le manoeuvre (146) est le
%  persan مزدور ; साईस le palefrenier (133) est l'arabe سائس, et
%  celui-la est ressorti — l'anglais en a fait « syce ». Restent
%  बढ़ई le charpentier, लोहार le forgeron, छवइया le couvreur et
%  सिखउआ l'apprenti, qui sont du vieux fonds. NEUF METIERS, CINQ
%  PROVENANCES, UN SEUL ECHAFAUDAGE.
%
%  LE PLAN (76) NE PEUT PAS S'APPELER « नकसा », PARCE QUE C'EST LE
%  NOM QUE CETTE COLONNE DONNE AUX PLANCHES. Le mot bhojpouri le plus
%  juste pour un plan d'architecte est नकसा, et c'est aussi le mot
%  choisi au cablage pour traduire « tabelo » — il est dans le titre
%  de chacun des seize fichiers. On ecrit donc « खाका ». UNE
%  TRADUCTION PEUT OCCUPER UN MOT DONT ELLE AURA BESOIN PLUS TARD, et
%  il n'y a rien a faire que de le constater et de reculer.
%
%  TRENTE-SIX PAIRES, ET CINQ MANQUEES : LE COMPTE EST LE MEME QU'AU
%  TABLEAU 3. Le tableau 4 en avait dix-huit et etait sorti net ; le
%  tableau 3 en avait vingt-deux et en avait manque six ; celui-ci en
%  a trente-six et en a manque cinq — le paratonnerre (40) derriere
%  le pignon (36), la scie (94) derriere le bois (69), l'echelle
%  (106) derriere le balcon (28), la ferrure (67) derriere
%  l'electricien (124), les marches (14) derriere le perron (9). LA
%  MESURE EST DONC CLAIRE : ce n'est pas le nombre de paires qui
%  fait la faute, c'est le moment ou la vigilance retombe. Le
%  tableau 4, plus court, avait ete verifie bloc a bloc de bout en
%  bout ; ici la verification a tenu jusqu'aux deux tiers.
%
%  ET LA MACRO \nl A ETE RECOPIEE DANS LA NOTE. Le bloc t05-noto-1
%  de l'ido coupe sa ligne au milieu de la liste des langues, et la
%  main a repris la coupe avec la macro. Une traduction ne suit
%  aucune ligne imprimee : kolonoj.py l'a prise pour ce qu'elle est,
%  une macro inconnue de cette colonne. C'est exactement la faute que
%  la colonne levantine avait faite a son tableau 2, et pour la meme
%  raison — on recopie l'apparat en meme temps que le texte.
%
%  TROIS DEGRES SUR LE PRONOM, ET LE TABLEAU EST UN DIALOGUE. C'est
%  le seul du livret ou l'on parle : trente-six attributions de
%  parole. Le bhojpouri distingue तू, तोहनी et रउआ, sur le pronom ET
%  sur le verbe, la ou l'ido n'a que « tu » et « vu ». Il faut donc
%  trancher a chaque replique, et c'est le texte qui tranche : les
%  adultes disent तू a Ioannes, Ioannes dit रउआ a son oncle, a
%  monsieur Blanc et a l'architecte. C'EST LA MEME CONTRAINTE QUE LE
%  « beau- » DU TABLEAU 4 : une langue plus fine que sa source oblige
%  a mieux lire la source. Ceci ne reprend rien a la colonne
%  javanaise : le javanais change de VOCABULAIRE entre ngoko et
%  krama, ce qui est autre chose ; ici c'est le pronom et l'accord.
% ===================================================================

%%K t05-apar-1 apar
\VUsaut{6.0mm}
\VUtitre{15.0pt}{60}{नकसा नं. ५}
\VUsaut{1.4mm}
\VUtitre{11.4pt}{0}{घर आ ओकर बनावट।}
\VUsaut{2.2mm}

%%K t05-tit-1 sub
\VUpk{10.2pt}{40}{नीक भेंट।}
\VUsaut{3.0mm}

%%K t05-01-1 p
\textsc{योहानेस}। --- चाचा, हमरा बहुते जाने के मन बा कि \VUgras{घर}
कइसे बनेला।

%%K t05-01-2 p
\textsc{चाचा} (80)। --- ई त बहुते आसान बा, हमार दोस्त, हमरा साथे आवऽ,
हम तोहरा ऊ देखाइब जवन बेर्नार्दुस बाबू (79) बनववत बाड़ें आ जवना में हम
रहब।

%%K t05-01-3 p
\textsc{योहानेस}। --- त एह घर के \VUgras{खाका} \textsuperscript{(76)}
के बनवले बा ?

%%K t05-01-4 p
\textsc{चाचा}। --- \VUgras{वास्तुकार} \textsuperscript{(75)} ; ऊ बहुते
साफ नक्सा देलें जवन मंजूर हो गइल, आ \VUgras{ठेकेदार}
\textsuperscript{(77)} काम सुरू कर देलें।

%%K t05-01-5 p
\textsc{योहानेस}। --- त ठेकेदार के हउवें ?

%%K t05-01-6 p
\textsc{चाचा}। --- ऊ ब्लांको बाबू हउवें, तोहार दोस्त याकोबुस के
बाबूजी। ऊ रूफो बाबू, वास्तुकार, के साथे मिल के मजूरन के काम करावेलें।

%%K t05-01-7 p
लऽ ! ठीक बेरा पर ब्लांको बाबू आपन \VUgras{पटरी} \textsuperscript{(78)}
लेके आ रूफो बाबू आपन खाका लेके आ गइलें।

%%K t05-01-8 p
राम-राम, बाबू लोग। रउआ सभ कइसन बानी ?

%%K t05-01-9 p
\textsc{रूफो बाबू}। --- बहुते नीक, धन्यवाद। राम-राम, योहानेस, ई बच्चा
त केतना बड़ हो गइल !

%%K t05-01-10 p
\textsc{चाचा}। --- हाँ, सचहूँ, ऊ छोट मरद हो गइल बा। ऊ बहुते गंभीर बा,
जवन कुछ ऊ देखेला ओकरा बारे में ओकरा के बतावे के परेला। आज ऊ जानल चाहत
बा कि घर कइसे बनेला।

%%K t05-tit-2 sub
\VUpk{10.2pt}{40}{मजूर लोग।}
\VUsaut{3.0mm}

%%K t05-01-11 p
\textsc{ब्लांको बाबू}। --- लऽ, हमरा साथे आवऽ, छोट दोस्त, हम तोहरा के
अभिए समझाइब, काहेकि हमरा ऊ लइका बहुते नीक लागेलें जे सीखल चाहेलें।

%%K t05-01-12 p
तू ओह मजूर के देखत बाड़ऽ जे हाथ में \VUgras{बेलचा} \textsuperscript{(82)}
पकड़ले बा : ऊ \VUgras{माटी खोदे वाला} \textsuperscript{(81)} ह। ऊ पानी के
पाइप धरे खातिर \VUgras{नाली} \textsuperscript{(46)} खोदत बा। ऊ ओह जगहा के
जमीन भी खोदले बा जहाँ घर बा, ताकि मिस्त्री चारो \VUgras{देवाल} उठा सके।
ओह देवाल के जवन हिस्सा जमीन के भीतर बा ओकरा के \VUgras{नेव}
\textsuperscript{(2)} कहल जाला। नेव के बीच के जगहा \VUgras{तहखाना}
\textsuperscript{(3)} बनावेला। ओह तहखाना में एगो छोट खिड़की होले जेकरा
\VUgras{तहखाना के झरोखा} \textsuperscript{(5)} कहल जाला, एगो
\VUgras{केवाड़} \textsuperscript{(4)} आ एगो सीढ़ी।

%%K t05-01-13 p
\VUgras{राजमिस्त्री} \textsuperscript{(108)} देवाल बनावत रहलें आ
\VUgras{केवाड़} \textsuperscript{(8)} आ \VUgras{खिड़की} \textsuperscript{(17)}
खातिर जगहा छोड़त रहलें।

%%K t05-01-14 p
\textsc{योहानेस}। --- बाकिर ऊ राजमिस्त्री त हमरा नइखे लउकत !

%%K t05-01-15 p
\textsc{ब्लांको बाबू}। --- ऊ अबहीं घर के भीतर के काम खतम कर लेलें। ऊ
\VUgras{घोड़सार} \textsuperscript{(54)} लगे बाड़ें, आपन \VUgras{मचान}
\textsuperscript{(110)} पर, आ \VUgras{धोबीघर} \textsuperscript{(56)} के
देवाल बनावत बाड़ें।

%%K t05-01-16 p
ओकरा लगे करनी, तसला आ \VUgras{साहुल} \textsuperscript{(109)} बा। बाकिर ई
नाँव तोहरा इयाद ना रहीहें।

%%K t05-01-17 p
\textsc{योहानेस}। --- जरूर रहीहें, काहेकि हम अभिए आपन चाचा से एगो छोट
करनी आ तसला माँगब, आ साहुल हम खुदे एगो डोरी से बना लेब।

%%K t05-tit-3 sub
\VUsaut{3.0mm}
\VUpk{10.2pt}{40}{सामान।}
\VUsaut{3.0mm}

%%K t05-01-18 p
\textsc{चाचा}। --- हा ! बहुते नीक। बाकिर हमरा बतावऽ, योहानेस, घर बनावे
खातिर का-का चाहीं ?

%%K t05-01-19 p
\textsc{योहानेस}। --- \VUgras{गढ़ल पाथर} \textsuperscript{(59)} आ
\VUgras{गारा} \textsuperscript{(65)} चाहीं।

%%K t05-01-20 p
\textsc{चाचा}। --- ठीक कहलऽ। ओह मनई के देखऽ, बड़का टोपी वाला, जे आपन
\VUgras{गाड़ी} \textsuperscript{(136)} पर बा। ऊ \VUgras{गाड़ीवान}
\textsuperscript{(135)} ह। ऊ रोज इहाँ \VUgras{पाथर के ढोंका}
\textsuperscript{(60)} ले आवेला। ऊ आपन गाड़ी अँगना में खाली करेला, आ
\VUgras{पाथर काटे वाला} \textsuperscript{(83)} ओह ढोंका के गढ़ेलें आ आपन
\VUgras{पाथर के आरी} \textsuperscript{(85)} से चीरेलें। \VUgras{मजूर}
\textsuperscript{(146)} ओकनी के राजमिस्त्री लगे ले जाला, जे ओकनी के
जगहा पर धरेलें।

%%K t05-01-21 p
ओही बेरा \VUgras{गारा सानेवाला} \textsuperscript{(87)} गारा तइयार करत
बा। ओकरा \VUgras{बालू} \textsuperscript{(62)}, \VUgras{चूना}
\textsuperscript{(63)} आ \VUgras{पानी} \textsuperscript{(64)} चाहीं। चूना
ओकरा लगे \VUgras{बोरा} \textsuperscript{(71)} में बा, \VUgras{मिस्त्री के
मददगार} \textsuperscript{(111)} ओकरा लगे \VUgras{ठेला} \textsuperscript{(112)}
पर बालू ले आवत बा आ \VUgras{सिखउआ} \textsuperscript{(113)} नल (58) लगे
बा। ऊ \VUgras{बाल्टी} \textsuperscript{(114)} भरत बा। गारा थोड़े देर में
तइयार हो जाई। \VUgras{गारा ढोवे वाला} \textsuperscript{(107)} ओकरा के
लेके सीढ़ी से मचान पर ऊपर ले जाला, जहाँ राजमिस्त्री घर के ओह ओर के काम
पूरा करत बाड़ें।

%%K t05-01-22 p
आ अबहीं बतावऽ, जे मजूर \VUgras{छत} \textsuperscript{(31)} के काम करेला ओकरा
के का कहल जाला ?

%%K t05-01-23 p
\textsc{योहानेस}। --- ऊ \VUgras{छवइया} \textsuperscript{(118)} ह।

%%K t05-tit-4 sub
\VUsaut{3.0mm}
\VUpk{10.2pt}{40}{घर के छत।}
\VUsaut{3.0mm}

%%K t05-01-24 p
\textsc{ब्लांको बाबू}। --- हाँ, बाकिर पहिले \VUgras{बढ़ई}
\textsuperscript{(115)} आवेला।

%%K t05-01-25 p
बढ़ई \VUgras{छत के ढाँचा} \textsuperscript{(35)} बनावेला। ऊ
\VUgras{कड़ी} \textsuperscript{(37)} के \VUgras{खूँटी} \textsuperscript{(117)}
से जोड़ेला। ऊ मजूर, ओहिजा ऊपर, आपन चाकर मखमली जाँघिया में, बढ़ई हउवें।
एगो छोट कड़ी पकड़ले बा आ दोसरका आपन \VUgras{टाँगी} \textsuperscript{(116)}
से खूँटी काटत बा। जे \VUgras{चिमनी} \textsuperscript{(41)} लगे बा ऊ
छवइया ह। ऊ (*) बल्ली पर पटरी ठोकेला, आ पटरी पर \VUgras{स्लेट}
\textsuperscript{(66)}। कबो-कबो ऊ खपड़ा धरेला।

%%K t05-noto-1 noto
\VUnotes{143.2mm}{%
(*) \VUgras{Balk-o} = जर्मन \textit{Balken}, अंगरेजी \textit{joist},
फ्रेंच \textit{solive}, इतालवी \textit{travicello}, रूसी
\textit{balka}, स्पेनी आ पुर्तगाली \textit{viga}.
ई तकनीकी सबद अबहीं ले लिहल ना गइल बा, बाकिर फ्रेंच \textit{solive} के
अरथ खातिर सुझावे लायक बा, जवन ना त trabo (फ्रेंच \textit{poutre}) ह ना
trabeto। ई चौकोर गढ़ल काठ के बल्ली ह जवन फर्स आ भितरी छत के थामेले।%
}

%%K t05-01-26 p
घोड़सार के छत \VUgras{खपड़ा वाला} \textsuperscript{(55)} बा, घर के छत
\VUgras{स्लेट वाला} \textsuperscript{(32)} बा आ बरामदा के छत
\VUgras{जस्ता वाला} \textsuperscript{(24)} बा।

%%K t05-01-27 p
\textsc{योहानेस}। --- का छवइया जस्ता के छत के भी काम करेला ?

%%K t05-01-28 p
\textsc{ब्लांको बाबू}। --- ना, ऊ काम \VUgras{जस्ता के मिस्त्री}
\textsuperscript{(121)} भा \VUgras{प्लंबर} करेला, ऊहे जे घर के छत पर
\VUgras{छत के खिड़की} \textsuperscript{(38)} बइठावेला। ऊ \VUgras{परनाला}
\textsuperscript{(43)} आ \VUgras{छत के नाली} \textsuperscript{(42)} भी
लगावेला। ऊ जस्ता आ टीन के जोड़ेला। ऊ \VUgras{बिजली-रोधक}
\textsuperscript{(40)} आ \VUgras{हवा के तीर} \textsuperscript{(39)}
\VUgras{छत के माथ} \textsuperscript{(36)} पर लगवले बा।

%%K t05-01-29 p
\textsc{योहानेस}। --- त घर बन गइल ?

%%K t05-tit-5 sub
\VUsaut{3.0mm}
\VUpk{10.2pt}{40}{औजार।}
\VUsaut{3.0mm}

%%K t05-01-30 p
\textsc{ब्लांको बाबू}। --- पूरा ना। \VUgras{पलस्तर वाला}
\textsuperscript{(99)} \VUgras{ईंटा} \textsuperscript{(61)} से ऊ भीत
बनावेला जवन घर के अलग-अलग कमरा में बाँटेला। ऊ \VUgras{पलस्तर}
\textsuperscript{(70)} से \VUgras{भितरी छत} \textsuperscript{(26)} आ
देवाल के लीपेला। अबहीं ऊ \VUgras{बरामदा} \textsuperscript{(33)} के भितरी
छत के काम करत बा। ओकरा लगे, राजमिस्त्री नियर, \VUgras{करनी}
\textsuperscript{(100)} आ \VUgras{तसला} \textsuperscript{(101)} बा।

%%K t05-01-31 p
ओकरा बाद काठ-मिस्त्री केवाड़, खिड़की, \VUgras{काठ के फर्स}
\textsuperscript{(16)} आ \VUgras{झिलमिली} \textsuperscript{(19)} बनावेला।

%%K t05-01-32 p
अबहीं ऊ अँगना में आपन \VUgras{काम के मेज} \textsuperscript{(90)} पर काम
करत बा। ओकरा लगे \VUgras{आरी} \textsuperscript{(94)} बा काठ (69) चीरे खातिर,
\VUgras{रंदा} \textsuperscript{(91)}, \VUgras{सिकंजा} \textsuperscript{(92)},
\VUgras{छेनी} \textsuperscript{(94 bis)}, \VUgras{परकार} \textsuperscript{(95)}
आ काठ के \VUgras{हथौड़ी} \textsuperscript{(95 bis)} बा।

%%K t05-01-33 p
\textsc{योहानेस}। --- हा ! बाकिर मिस्त्री के काम से बढ़ई के काम जादे
मजेदार बा। चाचा, रउआ हमरा खातिर काम के मेज, आरी आ रंदा खरीद दीं, काहेकि
हम जल्दिए मेदोर खातिर एगो झोपड़ी बनाइब।

%%K t05-01-34 p
\textsc{चाचा}। --- हाँ, बबुआ।

%%K t05-01-35 p
\textsc{योहानेस}। --- हम केतना खुस बानी ! आ ऊ के हउवें जे मुख्य केवाड़
लगे ठाढ़ बाड़ें ?

%%K t05-tit-6 sub
\VUsaut{3.0mm}
\VUpk{10.2pt}{40}{रंग-रोगन।}
\VUsaut{3.0mm}

%%K t05-01-36 p
\textsc{ब्लांको बाबू}। --- ऊ \VUgras{रंगसाज} \textsuperscript{(102)}
हउवें। ऊ आपन सीढ़ी पर \VUgras{रंग} \textsuperscript{(103)} के डिब्बा आ
आपन \VUgras{कूँची} \textsuperscript{(104)} लेके ठाढ़ बाड़ें। ऊ मुख्य
केवाड़ के काठ पर रंग लगावत बाड़ें ताकि ऊ जादे दिन चले।

%%K t05-01-37 p
ऊ कमरन में कागज भी चिपकावेलें।

%%K t05-01-38 p
\VUgras{परदा लगावे वाला} \textsuperscript{(107)}, जे \VUgras{सीढ़ी}
\textsuperscript{(106)} पर, \VUgras{बालकनी} \textsuperscript{(28)} पर बा,
\VUgras{परदा} \textsuperscript{(29)} लगावत बा।

%%K t05-01-39 p
जे मजूर बड़का खिड़की लगे ठाढ़ बा ऊ \VUgras{सीसा लगावे वाला}
\textsuperscript{(96)} ह। ऊ \VUgras{सीसा} \textsuperscript{(18)}
\VUgras{पुट्टी} \textsuperscript{(73)} से बइठावेला। ऊ दोसरका मजूर, जे
तहखाना के केवाड़ लगे घुटना टेकले बा, \VUgras{तालासाज}
\textsuperscript{(97)} ह। ऊ \VUgras{ताला} \textsuperscript{(6)} लगावत बा।
ओकर \VUgras{रेती} \textsuperscript{(98)} ओकरा लगे बा।

%%K t05-01-40 p
\textsc{योहानेस}। --- का ऊ मजूर भी तालासाज ह जे लोहा पर आपन
\VUgras{हथौड़ी} \textsuperscript{(84)} मारत बा ?

%%K t05-01-41 p
\textsc{ब्लांको बाबू}। --- ठीक-ठीक कहीं त ऊ लोहार (125) ह। ऊ
\VUgras{लोहा के सामान} \textsuperscript{(67)} गढ़त बा \VUgras{बिजली
मिस्त्री} \textsuperscript{(124)} खातिर, जे \VUgras{गाड़ीघर}
\textsuperscript{(51)} के छत पर बा। ओकरा लगे \VUgras{भट्ठी}
\textsuperscript{(126)}, \VUgras{निहाई} \textsuperscript{(127)} आ
\VUgras{संड़सी} \textsuperscript{(128)} बा।

%%K t05-01-42 p
बिजली मिस्त्री टेलीफोन के \VUgras{ताँबा} \textsuperscript{(44)} के
\VUgras{तार} बान्हत बा।

%%K t05-tit-7 sub
\VUpk{10.2pt}{40}{बगइचा के काम।}
\VUsaut{3.0mm}

%%K t05-01-43 p
\textsc{चाचा}। --- \VUgras{अँगना} \textsuperscript{(45)} के छोर पर,
\VUgras{लोहा के जाली} \textsuperscript{(47)} आ \VUgras{फाटक}
\textsuperscript{(48)} लगे, तू \VUgras{माली} \textsuperscript{(129)} के
देखत बाड़ऽ। ऊ \VUgras{छोट बगइचा} \textsuperscript{(49)} तइयार करत बा। ऊ
आपन \VUgras{कुदाल} \textsuperscript{(131)} से माटी गोड़ले बा, बीया बोवत
बा आ \VUgras{पंजा} \textsuperscript{(130)} से समतल करत बा। ओकरा बाद आपन
\VUgras{झारा} \textsuperscript{(132)} से ऊ \VUgras{क्यारी}
\textsuperscript{(150)} सींची आ पौधा बढ़ी।

%%K t05-01-44 p
\VUgras{साईस} \textsuperscript{(133)}, जे अभिए घोड़ा के मालिस कइले आ
ओकरा के दाना देले बा, \VUgras{गाड़ी} \textsuperscript{(52)} के स्पंज,
\VUgras{हाथ-पंप} \textsuperscript{(134)} आ बाल्टी भर पानी से साफ करत बा।
ओकरा बाद ऊ ओकरा के गाड़ीघर में लगा दी आ मोटका \VUgras{अरगल}
\textsuperscript{(53)} से केवाड़ बंद कर दी।

%%K t05-01-45 p
लऽ, अबहीं तू खुस बाड़ऽ, हमरा लागत बा ; तोहरा भरपूर समझ में आ गइल।

%%K t05-01-46 p
\textsc{योहानेस}। --- हाँ, चाचा, बाकिर हम खुसी से घर के भीतर भी देखतीं।

%%K t05-01-47 p
\textsc{चाचा}। --- ऊ काल्ह होई। रूफो बाबू तोहरा के खाका समझइहें आ हर
कमरा के नाँव बतइहें।

%%K t05-tit-8 sub
\VUsaut{3.0mm}
\VUpk{10.2pt}{40}{घर के भीतरी बनावट।}
\VUsaut{2.4mm}

%%K t05-apar-2 apar
{\centering\textit{(खाका देखीं।)}\par}
\VUsaut{2.4mm}

%%K t05-01-48 p
\textsc{वास्तुकार}। --- आवऽ, योहानेस, हम तोहरा के घर के अलग-अलग हिस्सा
अभिए घुमा देत बानी।

%%K t05-01-49 p
आवऽ, \VUgras{भुइयाँ-तल} \textsuperscript{(7)} में चलल जाव। ई रहल
\VUgras{घंटी} \textsuperscript{(11)} के बटन। हमनी के तीन गो \VUgras{डेग}
\textsuperscript{(14)} चढ़त बानी जा जवन \VUgras{चबूतरा}
\textsuperscript{(9)} के ह, \VUgras{देहरी} \textsuperscript{(10)} पार
करत बानी जा जवन \VUgras{मुख्य केवाड़} \textsuperscript{(8)} के ह, आ लऽ, हमनी के \VUgras{सीढ़ी}
\textsuperscript{(13)} के जड़ लगे बानी जा।

%%K t05-01-50 p
\textsc{योहानेस}। --- ई गलियारा ह।

%%K t05-01-51 p
\textsc{वास्तुकार}। --- ना, ई \VUgras{ड्योढ़ी} \textsuperscript{(12)} ह।
\VUgras{गलियारा} \textsuperscript{(a)} छोर पर बा। दाहिने ओर, ई रहल
\VUgras{बइठका} \textsuperscript{(b)} आ \VUgras{खाए के कमरा}
\textsuperscript{(c)} ; खाए के कमरा के सामने \VUgras{रसोई}
\textsuperscript{(d)} आ \VUgras{भंडार} \textsuperscript{(e)} बा, आ ओकरा
आगे \VUgras{लिखे के कमरा} \textsuperscript{(f)}। \VUgras{बरामदा}
\textsuperscript{(23)} आपन \VUgras{काँच के केवाड़} \textsuperscript{(25)}
के साथे बइठका लगे बा।

%%K t05-01-52 p
अबहीं हमनी के सीढ़ी चढ़ब जा। \VUgras{जँगला} \textsuperscript{(15)} पकड़ के
चलऽ आ डेग गिनऽ। चउदह गो बाड़ें। ई रहल \VUgras{सीढ़ी के चबूतरा}
\textsuperscript{(m)}, हमनी के पहिला \VUgras{मंजिल} \textsuperscript{(27)}
पर बानी जा।

%%K t05-tit-9 sub
\VUsaut{3.0mm}
\VUpk{10.2pt}{40}{पहिला मंजिल।}
\VUsaut{3.0mm}

%%K t05-01-53 p
सीढ़ी के सामने \VUgras{पाखाना} \textsuperscript{(20)} आ
\VUgras{सिंगार-कोठरी} \textsuperscript{(j)} बा। दाहिने ओर सुंदर
\VUgras{सुते के कमरा} \textsuperscript{(g)} बा, जवना के दू गो खिड़की घर
के \VUgras{अगवारी} \textsuperscript{(1)} पर बा। बाईं ओर \VUgras{नहाए के
कमरा} \textsuperscript{(1)} आ \VUgras{पाहुन के कमरा} \textsuperscript{(i)}
बा, जवना में बड़का \VUgras{कोठरी} \textsuperscript{(h)} बा। सुते के कमरा
लगे \VUgras{बच्चा के कमरा} \textsuperscript{(k)} बा। ऊ चाकर
\VUgras{छज्जा} \textsuperscript{(21)} से लागल बा। एह छज्जा के नीचे एगो
आउर पाखाना (20) बा, आ देवाल पर ई रहल \VUgras{घंट} \textsuperscript{(22)}
जवन खाना खातिर बाजेला।

%%K t05-01-54 p
\textsc{योहानेस}। --- चलीं \VUgras{दुसरका मंजिल} पर चलल जाव।

%%K t05-01-55 p
\textsc{वास्तुकार}। --- दुसरका मंजिल बा ही ना। छत के नीचे खाली
\VUgras{अटारी} \textsuperscript{(33)} आ कोठार (34) बा। हमनी के घुमाई
पूरा भइल, अबहीं नीचे उतरल जा सकेला।

%%K t05-01-56 p
\textsc{योहानेस}। --- धन्यवाद, बाबू, ई बहुते मजेदार रहल। हम जवन कुछ
देखनी ऊ सब अभिए आपन बाबूजी के सुनाइब।
